\documentclass[12pt]{article}
\usepackage[T1]{fontenc}
\usepackage[utf8]{inputenc}
\usepackage{lmodern}
\usepackage{textcomp}
\usepackage{enumitem}
\usepackage{geometry}
\usepackage{graphicx}
\usepackage{amsmath, amssymb}
\geometry{margin=1in}
\begin{document}
\begin{center}
Sociology - Undergraduate Level Assessment 15 \
Total Marks: 40 \
Answer all questions.
\end{center}

\section*{Multiple Choice (10 marks)}
\begin{enumerate}[label=\arabic*.]
\item The terms 'crisis of the 1970 s' is used to refer to:
\begin{enumerate}[label=(\alph*)]
    \item declining profits and rising unemployment
    \item the eradication of the welfare state
    \item rising divorce rates and the decline of the traditional family
    \item an unfortunate twist in fashion sensibility
\end{enumerate}
\item Sullivan's (2000) study suggested that the proportion of housework men did was greatest when:
\begin{enumerate}[label=(\alph*)]
    \item they had rediscovered themselves as 'new men'
    \item their wives were at home and nagged them all the time
    \item exciting gadgets like the hoover and electric iron were invented
    \item they were unemployed or both partners worked full time
\end{enumerate}
\item One of the difficulties in 'operationalizing' concepts like social class is that:
\begin{enumerate}[label=(\alph*)]
    \item definitions and indicators can vary, making valid comparisons problematic
    \item there are no reliable indicators of such widely contested ideas
    \item it takes all the fun out of armchair theorising
    \item it has little use for applied, empirical research about the topic
\end{enumerate}
\item In idealized views of science, the experimental method is said to involve:
\begin{enumerate}[label=(\alph*)]
    \item testing out new research methods to see which one works best
    \item isolating and measuring the effect of one variable upon another
    \item using personal beliefs and values to decide what to study
    \item interpreting data subjectively, drawing on theoretical paradigms
\end{enumerate}
\item White-collar crime is low in visibility because:
\begin{enumerate}[label=(\alph*)]
    \item it involves only small amounts of money
    \item the proletariat can outsmart the bourgeoisie
    \item the police turn a blind eye to corporate crime
    \item it goes undetected in the context of everyday business transactions
\end{enumerate}
\end{enumerate}

\section*{True / False (10 marks)}
\textit{Answer True or False}
\begin{enumerate}[label=\arabic*.]
\setcounter{enumi}{5}
\item True or False: 'Eyeballing' refers to scanning a table to identify general patterns and significant figures in data.
\item True or False: Terrorism is characterized by its organized nature and often global scale, distinct from typical crimes described by the Chicago School.
\item True or False: The demographic transition involves a decline in birth rates, increased life expectancy, and an ageing population.
\item True or False: Extensive paperwork is considered a fundamental aspect of Weber's ideal type of bureaucracy.
\item True or False: Marx believed that religion would diminish following a socialist revolution that eliminated capitalist ideology.
\end{enumerate}

\section*{Long Form Response (20 marks)}
\textit{Answer each question with a clear and complete explanation.}
\begin{enumerate}[label=\arabic*.]
\setcounter{enumi}{10}
\item Discuss the process of socialization and analyze its significance in shaping individual identity and facilitating membership in society.
\item Discuss the concept of vertical mobility in sociology, analyzing its implications on social structure and individual identity within various societal contexts.
\end{enumerate}

\vfill
\noindent\textit{End of Paper}
\end{document}